
\documentclass[11pt,a4paper,oneside]{article}
\usepackage[top=0.85in,bottom=0.85in,left=0.95in,right=0.95in,headheight=14pt]{geometry}
\usepackage[T1]{fontenc}
\usepackage{lmodern}
\usepackage{microtype}
\usepackage{amsmath,amssymb,amsthm,mathtools,bm}
\usepackage{booktabs,array,tabularx,longtable}
\usepackage{xcolor}
\usepackage{enumitem}
\usepackage{fancyhdr}
\usepackage{titlesec}
\usepackage[colorlinks=true,linkcolor=black,urlcolor=blue,citecolor=black]{hyperref}

\setlist{leftmargin=1.4em,itemsep=0.25em,topsep=0.25em}
\titleformat{\section}{\Large\bfseries}{\S\thesection}{0.8em}{}
\titleformat{\subsection}{\large\bfseries}{\thesubsection}{0.8em}{}
\titleformat{\subsubsection}{\normalsize\bfseries}{\thesubsubsection}{0.8em}{}

\newtheorem{definition}{Definition}[section]
\newtheorem{axiom}{Axiom}[section]
\newtheorem{proposition}{Proposition}[section]
\newtheorem{remark}{Remark}[section]

%% ================================================================
%%  MACRO SET --- EVM-V v1.0
%%  Harmonised with IKK II v2.0, IKK III v2.0, IKK IV v1.0
%% ================================================================

%% ---- Identity objects (IKK registry) ----
\newcommand{\Ia}{I_a}
\newcommand{\Itrue}{\mathfrak{I}_{\mathrm{true}}}
\newcommand{\Ihat}{\widehat{I}_{a|b}}

%% ---- Typed spaces ----
\newcommand{\Osp}{\mathcal{O}_b}
\newcommand{\Pab}{\mathcal{P}_{a\to b}}
\newcommand{\Rresid}{\mathfrak{R}_{a|b}}

%% ---- Observed accelerations ----
\newcommand{\aobs}{a_{\mathrm{obs}}}
\newcommand{\abar}{a_{\mathrm{bar}}}
\newcommand{\aPi}{a_{\Pi}}
\newcommand{\aMOND}{a_{\mathrm{MOND}}}
\newcommand{\aPiM}{a_{\Pi,\mathrm{MOND}}}
\newcommand{\ascale}{a_s}

%% ---- Velocities ----
\newcommand{\vobs}{v_{\mathrm{obs}}}
\newcommand{\vbar}{v_{\mathrm{bar}}}
\newcommand{\vPi}{v_{\Pi}}
\newcommand{\vmod}{v_{\mathrm{model}}}
\newcommand{\vgas}{v_{\mathrm{gas}}}
\newcommand{\vdisk}{v_{\mathrm{disk}}}
\newcommand{\vbulge}{v_{\mathrm{bulge}}}
\newcommand{\vcirc}{v_{\mathrm{circ}}}
\newcommand{\vN}{v_N}
\newcommand{\vfive}{v_5}

%% ---- Model parameters ----
\newcommand{\LamPi}{\Lambda_{\Pi}}
\newcommand{\pzero}{p_0}
\newcommand{\rc}{r_c}
\newcommand{\rzero}{r_0}

%% ---- Hidden momentum ----
\newcommand{\pu}{p_u}
\newcommand{\pw}{p_w}      % deprecated; kept for Looking Up cross-references only

%% ---- Scale coordinate ----
\newcommand{\usc}{u}                    % log-scale chart u = ln(l/l_0)
\newcommand{\ellz}{\ell_0}             % reference scale

%% ---- Projection efficiency ----
\newcommand{\etaPi}{\eta_{\Pi}}
\newcommand{\etaab}{\eta_{ab}}

%% ---- Gate function ----
\newcommand{\GPi}{G_{\Pi}}

%% ---- Mass-to-light ratios ----
\newcommand{\Upsdisk}{\Upsilon_{\mathrm{disk}}}
\newcommand{\Upsbulge}{\Upsilon_{\mathrm{bulge}}}

%% ---- Baryonic densities ----
\newcommand{\Sbar}{\Sigma_{\mathrm{bar}}}
\newcommand{\rhobar}{\rho_{\mathrm{bar}}}
\newcommand{\rhoNFW}{\rho_{\mathrm{NFW}}}

%% ---- Galaxy structure ----
\newcommand{\Mstar}{M_{\star}}
\newcommand{\fgas}{f_{\mathrm{gas}}}
\newcommand{\Rd}{R_d}

%% ---- Standard shorthands ----
\newcommand{\Oh}{\mathcal{O}}
\newcommand{\RR}{\mathbb{R}}
\newcommand{\CC}{\mathbb{C}}

%% ================================================================

\pagestyle{fancy}
\fancyhf{}
\fancyhead[L]{\small\textit{Inverse Kaluza-Klein Framework --- EVM Paper Set}}
\fancyhead[R]{\small\textit{EVM-V v1.0}}
\fancyfoot[C]{\thepage}
\renewcommand{\headrulewidth}{0.4pt}
\setlength{\parskip}{0.55em}
\setlength{\parindent}{0pt}

\begin{document}

\begin{titlepage}
\centering
\vspace*{1.1cm}
{\large\scshape VSEPR-SIM Theoretical Division / EVM Paper Set\par}
\vspace{1.1cm}
{\Huge\bfseries EVM-V\par}
\vspace{0.35cm}
{\Large\bfseries Galactic Rotation Residuals as Projection-Access Signals\par}
\vspace{0.2cm}
{\large Version 1.0\par}
\vspace{1.5cm}

\begin{abstract}
\noindent
This paper develops EVM-V, the first large-scale observational evidence-mapping
document for the Inverse Kaluza-Klein Scale Symmetry framework. The preceding
\emph{Looking Up} paper proposed that galactic rotation anomalies may be interpreted
not as direct evidence for unseen matter alone, nor as a purely empirical modification
of Newtonian dynamics, but as a projection-residual effect associated with inaccessible
scale-depth structure. In that formulation, the observed radial acceleration contains a
baryonic Newtonian term plus an additional fifth-force-like contribution sourced by
hidden-axis momentum and activated only at galactic scales.

The present paper does not claim discovery of a new force. Instead, it defines a
reproducible data-analysis program for testing whether galactic residual accelerations
can be mapped onto a projection-access model. The central empirical object is
\[
    \aPi(r) = \aobs(r) - \abar(r),
\]
where $\aobs = \vobs^2/r$ is the observed centripetal acceleration and
$\abar = \vbar^2/r$ is the acceleration predicted from baryonic matter alone.
The paper treats the SPARC database as the primary benchmark, which contains 175
disk galaxies with \emph{Spitzer} 3.6\,$\mu$m photometry, high-quality
H\,\textsc{i}/H$\alpha$ rotation curves, and mass models suitable for baryonic
acceleration reconstruction.

EVM-V compares three interpretations of the residual: conventional dark matter halo
fitting, MOND-like phenomenological acceleration laws, and IKK projection-residual
dynamics. The immediate goal is not to outperform $\Lambda$CDM or MOND, but to
determine whether the projection-residual model yields stable population-level
parameters, physically meaningful correlations, and non-arbitrary failure conditions.
A model is considered scientifically useful only if its universal parameters remain
stable across galaxies and its residual structure cannot be reduced to a disguised halo
fit.

\medskip
\noindent\textbf{Keywords:}
galactic rotation curves; SPARC; radial acceleration relation; MOND; dark matter;
fifth force; projection residual; Inverse Kaluza-Klein; scale symmetry; baryonic
acceleration; hidden-axis momentum; effective force law.
\end{abstract}
\end{titlepage}

%% ================================================================
\section{Introduction}
%% ================================================================

The rotation-curve problem is one of the cleanest observational failures of naive
baryonic Newtonian dynamics. In spiral and irregular disk galaxies, observed orbital
velocities often remain approximately flat at large radii instead of declining as
expected from visible mass alone. The standard interpretation is that galaxies are
embedded in extended dark matter halos. This interpretation is mathematically productive
and cosmologically powerful, but the underlying non-baryonic particle has not been
directly detected. Modified gravity approaches, including MOND, instead treat the
discrepancy as a modification of the acceleration law at low acceleration. The SPARC
database and the radial acceleration relation have made this comparison unusually sharp
because they connect observed acceleration and baryonic acceleration across thousands of
resolved galactic radii.

The earlier \emph{Looking Up} paper proposed a third interpretive route. Instead of
treating the missing acceleration as either invisible mass or a primitive modification
to gravity, it framed the anomaly as an unresolved projection residual. In the language
of IKK, ordinary observation does not recover full physical identity:
\[
    \Osp^{(a)} = \Pab[\Ia], \qquad \Osp^{(a)} \neq \Ia.
\]
The observable object is a projection, not the primitive identity. The IKK notation
registry explicitly preserves this distinction through the canonical rule
\[
    \Itrue \;\neq\; I_{\mathrm{top}} \;\neq\; \Osp^{(a)}.
\]
This distinction matters because a residual in galactic dynamics can be treated not as
a free excuse term, but as the mismatch between gravitationally inferred structure and
electromagnetically resolved structure.

The prior \emph{Looking Up} model introduced hidden-axis momentum
\[
    p_w^{(i)} = m_i\,\frac{dw_i}{d\tau}
\]
as the effective fifth-force charge, and proposed a logarithmic radial acceleration
contribution which behaves as a $1/r$ acceleration at galactic radii, contributing an
approximately constant term to $v^2(r)$.

EVM-V reframes this result as a data-analysis protocol, not a declaration of victory.
Its central question is:
\begin{quote}
\textit{Can galactic acceleration residuals be mapped onto stable projection-access
dynamics?}
\end{quote}
That is a useful question. It does not require pretending dark matter is dead, MOND is
complete, or the fifth force has descended from the heavens carrying a grant proposal.

%% ================================================================
\section{Relation to \emph{Looking Up}}
%% ================================================================

The \emph{Looking Up} paper developed a fifth-dimensional momentum framework for
galactic rotation without invoking particle dark matter. Its abstract framed dark-sector
signatures as ``gravitationally visible but electromagnetically unresolved identity
content,'' i.e.\ a projection gap rather than invisible matter.

The present paper differs in purpose.

\emph{Looking Up} proposed a mechanism:
\[
    v^2(r) = \vN^2(r) + \vfive^2(r).
\]
EVM-V proposes an evidence pipeline:
\[
    \text{data}
    \;\to\; \aobs
    \;\to\; \abar
    \;\to\; \aPi
    \;\to\; \text{model comparison}
    \;\to\; \text{parameter stability}
    \;\to\; \text{failure conditions.}
\]
The paper therefore has a narrower and more scientific burden. It asks whether the
fifth-force-like term from \emph{Looking Up} can survive direct comparison against
SPARC rotation curves, the radial acceleration relation, MOND-like scaling, and
conventional dark matter halo fits.

The earlier paper already identified the correct primary benchmark: SPARC
rotation-curve fitting, radial acceleration relation analysis, compact-system screening
tests, population tests, and morphological correlations. EVM-V expands those into a
full research design.

%% ================================================================
\section{Notation and Framework Alignment}
%% ================================================================

The notation used here follows the IKK registry where possible.

\subsection{Identity and projection}

Let $\Ia$ denote primitive or source identity at scale $S_a$. Observation at scale
$S_b$ produces
\[
    \Osp^{(a)} = \Pab[\Ia],
\]
where $\Pab$ is the projection operator from source scale $S_a$ into observable scale
$S_b$. The registry gives the canonical operational form as
\[
    \Pab = \mathcal{M}_b \circ \mathcal{R}_b \circ \mathcal{C}_{ab},
\]
or, in simple form,
\[
    \Pab = \mathcal{R}_b \circ \mathcal{C}_{ab},
\]
where $\mathcal{C}_{ab}$ is the cross-scale coupling rule and $\mathcal{R}_b$ is the
resolution operator of observing scale $S_b$.

\subsection{Scale coordinate}

The legacy coordinate $w$ appears in \emph{Looking Up} because that paper was written
before the current notation registry stabilised. In this paper, the preferred coordinate
is
\[
    \usc = \ln\!\left(\frac{\ell}{\ellz}\right),
\]
where $\ell$ is the relevant physical scale and $\ellz$ is a reference scale. The
registry marks $\usc$ as the accepted logarithmic scale coordinate and marks $w$ as
deprecated in new writing.

To preserve continuity with \emph{Looking Up}, we write $\pu$ as the modernised analogue
of hidden-axis momentum, with the understanding that it descends from the earlier $\pw$
construction:
\[
    \pw \;\longrightarrow\; \pu.
\]
The replacement is conceptual as well as symbolic. The fifth direction is no longer
treated as an ordinary hidden spatial coordinate. It is a scale-access coordinate
governing how identity content becomes dynamically visible or invisible under projection.

\subsection{Residual acceleration}

The observable acceleration residual is defined as
\[
    \aPi(r) = \aobs(r) - \abar(r).
\]
The symbol $\Pi$ is used here only as a typed residual label, not as a generic
projection operator. The notation registry explicitly warns against bare $\Pi$ because
of projection and chemistry collisions. Tiny symbol hygiene, gigantic reduction in
future suffering.

%% ================================================================
\section{Physical Hypothesis}
%% ================================================================

The working hypothesis of EVM-V is:
\[
    \aobs(r) = \abar(r) + \aPi(r,\usc),
\]
where $\aPi(r,\usc)$ is an effective projection-residual acceleration.

This is not initially claimed to be a fundamental force in the Standard Model sense.
It is a fifth-force-like term in the limited phenomenological sense that it contributes
an additional acceleration:
\[
    F_{\mathrm{eff}} = m\,\aobs = m\,\abar + m\,\aPi.
\]
The interpretation is:
\begin{quote}
\textit{The residual is observable; the ontology of the residual is inferred.}
\end{quote}
This is the central scientific posture of the paper. The residual acceleration can be
measured from rotation-curve and baryonic-mass data. The claim that it reflects
projection-access structure is a model interpretation, and that interpretation must earn
credibility through parameter stability, cross-galaxy consistency, and comparison with
existing alternatives.

%% ================================================================
\section{Empirical Baseline: Rotation Curves and the RAR}
%% ================================================================

\subsection{Observed acceleration}

For each galaxy, the observed centripetal acceleration at radius $r_j$ is
\[
    \aobs(r_j) = \frac{\vobs^2(r_j)}{r_j}.
\]
The observational inputs are radius, observed velocity, and velocity uncertainty:
\[
    \bigl\{r_j,\;\vobs(r_j),\;\sigma_v(r_j)\bigr\}_{j=1}^{N}.
\]

\subsection{Baryonic acceleration}

The baryonic velocity contribution is decomposed as
\[
    \vbar^2(r) = \vgas^2(r) + \Upsdisk\,\vdisk^2(r) + \Upsbulge\,\vbulge^2(r),
\]
where $\Upsdisk$ and $\Upsbulge$ are stellar mass-to-light ratios. The baryonic
acceleration is
\[
    \abar(r) = \frac{\vbar^2(r)}{r}.
\]
This is why SPARC is the natural first benchmark: it provides rotation curves and
baryonic mass models for 175 disk galaxies across wide ranges of stellar mass, surface
brightness, and gas fraction.

\subsection{Radial acceleration relation}

The empirical radial acceleration relation connects $\aobs$ and $\abar$ across resolved
radii in disk galaxies. McGaugh, Lelli, and Schombert reported a tight relation using
SPARC data, making it one of the most direct targets for any projection-residual theory.

The EVM-V model must therefore be tested against both individual rotation curves and
the population-level relation
\[
    \aobs = f(\abar).
\]
A model that fits individual galaxies but destroys the radial acceleration relation is
not a success. It is just a mathematical ferret loose in the data.

%% ================================================================
\section{MOND as Phenomenological Baseline}
%% ================================================================

MOND provides the most relevant phenomenological comparison because it already treats
the anomaly as an acceleration-scale effect. In the deep-MOND regime,
\[
    \aMOND \approx \sqrt{a_0\,\abar},
\]
where $a_0$ is the characteristic acceleration scale.

The corresponding MOND residual is
\[
    \aPiM(r) = \aMOND(r) - \abar(r).
\]
EVM-V uses this not as an assumption, but as a benchmark. The IKK projection-residual
model should reproduce MOND-like behaviour in the regime where MOND is empirically
successful:
\[
    \aPi(r,\usc) \approx \aPiM(r) \quad \text{when } \abar \ll a_0.
\]
The interpretation differs:
\begin{quote}
MOND describes the low-acceleration law.\\[2pt]
IKK attempts to explain \emph{why} a scale-access law appears.
\end{quote}
This is the correct relationship. EVM-V is not ``MOND but with extra mist.'' It is a
proposed ontological mapping beneath MOND-like phenomenology. Whether that mapping is
worth anything depends on the analysis.

Li et al.\ fitted the radial acceleration relation to individual SPARC galaxies and
found that allowing galaxy-to-galaxy variation in the acceleration scale did not improve
the fits, which strengthens the burden on any IKK-like model: any additional
scale-access parameter must either reduce to a nearly universal value or be
independently tied to observable structure.

%% ================================================================
\section{Projection-Residual Model}
%% ================================================================

\subsection{Minimal residual form}

The minimal EVM-V acceleration model is
\[
    \aobs(r) = \abar(r) + \aPi(r,\usc).
\]
A direct descendant of \emph{Looking Up} is
\[
    \aPi(r,\usc) = \LamPi\,\ln\!\left(1 + \frac{|\pu(r)|}{\pzero}\right)
    \frac{\GPi(r)}{r + \rc},
\]
where
\[
    \GPi(r) = \frac{1}{1 + e^{-k(r-\rzero)}}.
\]
This gives the full form:
\[
    \aPi(r,\usc) = \LamPi\,\ln\!\left(1 + \frac{|\pu(r)|}{\pzero}\right)
    \cdot\frac{1}{1 + e^{-k(r-\rzero)}}
    \cdot\frac{1}{r+\rc}.
\]
The corresponding velocity contribution is
\[
    \vPi^2(r) = r\,\aPi(r,\usc),
\]
so that
\[
    \vmod^2(r) = \vbar^2(r) + \vPi^2(r).
\]
This is the modern-notation version of the \emph{Looking Up} rotation-curve equation,
where the fifth-force velocity term was written with $\Lambda_5$, $\pw$, $\pzero$,
$\rzero$, $k$, and $\rc$.

\subsection{Outer-disk limit}

At large galactic radii, $r \gg \rzero$ and $r \gg \rc$, so $\GPi(r) \to 1$ and
$r/(r+\rc) \to 1$. Then
\[
    \vPi^2(r) \approx \LamPi\,\ln\!\left(1 + \frac{|\pu(r)|}{\pzero}\right).
\]
If $\pu(r)$ varies slowly across the outer disk, $\vPi^2$ is approximately constant,
producing flat rotation curves. This is exactly the key result of the original
\emph{Looking Up} derivation.

\subsection{Projection-access interpretation}

The term $\pu(r)$ is not directly observed. It is a proxy for scale-access momentum
or hidden projection activity. Formally,
\[
    \pu(r) \sim \frac{\partial \Ia}{\partial \usc},
\]
meaning it measures how strongly the identity content changes along the scale-access
coordinate. This should not be oversold. It is a modelling bridge, not a measured
microscopic quantity.

Define a projection-access efficiency $\etaPi(r,\usc) \in [0,1]$, and write the
residual more generally as
\[
    \aPi(r,\usc) = \etaPi(r,\usc)\,\ascale(r),
\]
where $\ascale(r)$ is the scale-residual acceleration kernel. The logarithmic
fifth-momentum model is then one candidate specification:
\[
    \etaPi(r,\usc)\,\ascale(r)
    = \LamPi\,\ln\!\left(1+\frac{|\pu(r)|}{\pzero}\right)
      \GPi(r)\,\frac{1}{r+\rc}.
\]
The notation registry currently treats $\etaab$ as a provisional projection efficiency
or coupling strength, which fits this use but requires explicit definition in this paper.

%% ================================================================
\section{Hidden-Momentum Proxy Models}
%% ================================================================

The single largest modelling risk is $\pu(r)$. If it is freely chosen per galaxy, the
model degenerates into a disguised halo fit. The prior \emph{Looking Up} paper already
identified this as the hidden-momentum proxy problem.

EVM-V therefore defines four proxy tiers.

\subsubsection*{Tier I: Circular-kinematic proxy}
\[
    |\pu(r)| = \alpha_u\,\Mstar\,\vcirc(r).
\]
Simplest model. Ties scale-access momentum to local circular motion. Risks circularity
because $\vcirc$ is close to the observable being fitted. Use only as a baseline
sanity check.

\subsubsection*{Tier II: Surface-density proxy}
\[
    |\pu(r)| = \alpha_u\,\Sbar(r)\,\vbar(r).
\]
Here $\Sbar(r)$ is the baryonic surface density. Ties the residual to local baryonic
structure. This is the preferred first serious proxy.

\subsubsection*{Tier III: Momentum-flux proxy}
\[
    |\pu(r)| = \alpha_u\,\rhobar(r)\,\vbar(r).
\]
Treats hidden-axis sourcing as related to local baryonic momentum density. More
physically motivated but requires assumptions about vertical disk structure and
density reconstruction.

\subsubsection*{Tier IV: Morphology-conditioned proxy}
\[
    |\pu(r)| = \alpha_u\,f\!\left(\Sbar,\,\Mstar,\,\fgas,\,\Rd,\,\mathrm{morph.}\right)
    \vbar(r).
\]
Should be used only after Tiers I--III are evaluated. Otherwise the model will do what
models love doing when unsupervised: eat degrees of freedom and call it insight.

%% ================================================================
\section{Data Sources}
%% ================================================================

\subsection{Primary dataset: SPARC}

The primary dataset is SPARC, which contains 175 nearby late-type galaxies with
\emph{Spitzer} 3.6\,$\mu$m photometry and high-quality H\,\textsc{i}/H$\alpha$
rotation curves. Required SPARC fields per data point:
\[
    r_j,\;\vobs(r_j),\;\sigma_v(r_j),\;\vgas(r_j),\;\vdisk(r_j),\;\vbulge(r_j),
    \;L_{3.6},\;D,\;i.
\]
Here $D$ is distance and $i$ is inclination.

\subsection{Derived quantities}

For each data point:
\begin{align}
    \aobs(r_j) &= \frac{\vobs^2(r_j)}{r_j}, \\[4pt]
    \vbar^2(r_j) &= \vgas^2(r_j)
        + \Upsdisk\,\vdisk^2(r_j) + \Upsbulge\,\vbulge^2(r_j), \\[4pt]
    \abar(r_j) &= \frac{\vbar^2(r_j)}{r_j}, \\[4pt]
    \aPi^{\,\mathrm{data}}(r_j) &= \aobs(r_j) - \abar(r_j).
\end{align}

\subsection{Velocity uncertainty propagation}

Velocity uncertainty propagates to acceleration as
\[
    \sigma_{\aobs} = \frac{2\vobs\,\sigma_v}{r}.
\]
Uncertainty from stellar mass-to-light ratios $\Upsdisk$, $\Upsbulge$ must also be
included as systematic floors on $\abar$.

\subsection{Future dataset expansion}

BIG-SPARC has been proposed as a much larger successor database, targeting approximately
4000 galaxies using H\,\textsc{i} data cubes from public telescope archives and
near-infrared photometry from WISE. EVM-V begins with SPARC because it is established
and manageable. BIG-SPARC belongs in EVM-VI or a later population-scale validation
paper, assuming the model survives initial contact with reality. Harsh environment,
reality.

%% ================================================================
\section{Analysis Pipeline}
%% ================================================================

\subsection{Preprocessing}

For each galaxy:
\begin{enumerate}
  \item Load SPARC rotation-curve and mass-model data.
  \item Remove flagged or low-quality points according to SPARC quality metadata.
  \item Convert radii to SI units or use consistent astrophysical units
        (kpc, km\,s$^{-1}$, m\,s$^{-2}$).
  \item Compute $\vbar$, $\aobs$, $\abar$, and $\aPi^{\mathrm{data}}$.
  \item Propagate velocity uncertainties to $\sigma_{\aobs}$.
  \item Include uncertainty from mass-to-light ratios $\Upsdisk$, $\Upsbulge$ as
        systematic floors on $\abar$.
\end{enumerate}

\subsection{Baseline residual extraction}

Compute
\[
    \Delta v^2(r) = \vobs^2(r) - \vbar^2(r)
\]
and
\[
    \aPi^{\mathrm{data}}(r) = \frac{\Delta v^2(r)}{r}.
\]
This gives the empirical residual without assuming IKK, MOND, or halo structure.

\subsection{Model fitting}

Fit
\[
    \vmod^2(r;\,\Theta) = \vbar^2(r)
    + \LamPi\,r\,\ln\!\left(1+\frac{|\pu(r)|}{\pzero}\right)
      \GPi(r)\,\frac{1}{r+\rc}.
\]
Parameter vector:
\[
    \Theta = \{\LamPi,\;\pzero,\;\rc,\;\rzero,\;k,\;\alpha_u\}.
\]
Depending on proxy tier, $\alpha_u$ may be universal, galaxy-class-dependent, or
absorbed into $\pzero$. Minimise
\[
    \chi^2(\Theta) = \sum_{j=1}^{N}
    \left[\frac{\vobs(r_j) - \vmod(r_j;\,\Theta)}{\sigma_v(r_j)}\right]^2.
\]
Also compute acceleration-space loss:
\[
    \chi^2_a(\Theta) = \sum_{j=1}^{N}
    \left[\frac{\aobs(r_j) - a_{\mathrm{model}}(r_j;\,\Theta)}{\sigma_{a,j}}\right]^2.
\]
Both should be reported because velocity-space and acceleration-space fitting can weight
data differently.

\subsection{Comparison models}

For each galaxy, compare against:
\begin{enumerate}
  \item \textbf{Newtonian baryons only}: $\vmod^2 = \vbar^2$.
  \item \textbf{MOND/RAR}: $\aMOND = \nu({\abar}/{a_0})\,\abar$, with a chosen
        interpolation function $\nu$.
  \item \textbf{Dark matter halo}: At minimum NFW,
        $\rhoNFW(r) = \rho_s\bigl[(r/r_s)(1+r/r_s)^2\bigr]^{-1}$;
        optionally Burkert or Einasto profiles.
  \item \textbf{IKK projection-residual}: $\vmod^2 = \vbar^2 + \vPi^2$.
\end{enumerate}

\subsection{Model-selection metrics}

Report $\chi^2_\nu$, AIC, BIC, $\Delta$AIC, $\Delta$BIC, where
\[
    \mathrm{AIC} = 2k - 2\ln\mathcal{L}, \qquad
    \mathrm{BIC} = k\ln N - 2\ln\mathcal{L}.
\]
This matters because the projection model must not win merely by having more knobs.
Humanity already invented epicycles once. We do not need a premium subscription version.

%% ================================================================
\section{Population-Level Tests}
%% ================================================================

A galaxy-by-galaxy fit is insufficient. The decisive test is whether the model has
stable population-level structure.

\subsection{Universal parameter stability}

The parameters $\LamPi$, $\pzero$, $\rc$ should be approximately universal or should
vary according to a physically justified predictor. Define the population stability
statistic:
\[
    S_\theta = \frac{\mathrm{std}[\log_{10}\theta_g]}
                    {\sigma_{\theta,\mathrm{allowed}}},
\]
where $\theta_g$ is the fitted value for galaxy $g$. A parameter passes if
$S_\theta \leq 1$. A practical criterion is
\[
    \max_g(\theta_g)\,/\,\min_g(\theta_g) < 3
    \quad\text{(strong stability)}
\]
and $< 10$ for weak admissibility. Beyond that, the model starts wearing a fake halo
moustache.

\subsection{Correlation with galaxy properties}

If parameters vary, variation must correlate with observables: $\Mstar$,
$M_{\mathrm{gas}}$, $\fgas$, $\Rd$, $\Sigma_0$, surface brightness, morphology. For
each parameter $\theta$, test:
\[
    \log\theta = \beta_0 + \beta_1\log\Mstar + \beta_2\log\fgas
    + \beta_3\log\Rd + \epsilon.
\]
No correlation plus high variation means failure.

\subsection{RAR preservation}

The projection-residual model must preserve the observed tightness of the RAR. Define
residuals:
\[
    \delta_{\mathrm{RAR}} = \log\aobs - \log a_{\mathrm{model}}.
\]
Report $\sigma_{\mathrm{RAR}} = \mathrm{std}(\delta_{\mathrm{RAR}})$. A serious IKK
model should approach the scatter achieved by empirical RAR/MOND-style fits. Li et al.\
found small residual scatter when fitting the RAR to SPARC galaxies, so EVM-V must
treat that as a high bar, not a decorative citation.

%% ================================================================
\section{Screening and Compact-System Consistency}
%% ================================================================

The original \emph{Looking Up} model uses a radial sigmoid gate
\[
    \GPi(r) = \frac{1}{1 + e^{-k(r-\rzero)}}
\]
to suppress the fifth-force-like contribution below galactic scales. EVM-V must turn
this into a testable consistency requirement.

\subsection{Solar-system constraint}

For solar-system radii, $r \ll \rzero$, so $\GPi(r) \approx 0$. The model must satisfy
\[
    |\aPi(r_{\mathrm{solar}})| \ll a_{\mathrm{bound}},
\]
where $a_{\mathrm{bound}}$ is the tightest available experimental bound on anomalous
accelerations at solar-system scales.

\subsection{Wide-binary and dwarf-galaxy transition zone}

The danger zone is not molecular dynamics. The danger zone is intermediate gravitational
systems: $r \sim 10^{-3}$ to $10^1$\,kpc. The model must specify whether the
transition is determined by absolute radius, system mass, acceleration scale,
environmental density, or scale coordinate $\usc$.

A pure radius gate may be too crude. A more defensible gate is
$\GPi = \GPi(\usc,\,\abar,\,\Sbar)$ rather than $\GPi = \GPi(r)$ alone. A
radius-only gate is mathematically convenient but physically suspicious. Nature rarely
respects our desire for easy logistic functions, rude little universe that it is.

%% ================================================================
\section{Cluster-Scale Consistency}
%% ================================================================

Galaxy clusters are the hardest test for modified-gravity and fifth-force-like
frameworks. The \emph{Looking Up} paper explicitly notes that the Bullet Cluster creates
a major constraint because the lensing mass follows the collisionless galaxies rather
than the X-ray gas, even though the gas contains much of the baryonic mass.

EVM-V therefore requires a separate cluster analysis. The model must explain why
$\aPi^{\mathrm{cluster}}$, or the equivalent lensing residual, follows collisionless
stellar components rather than collisional gas.

Possible mechanisms include:
\begin{itemize}
  \item $\pu$ is sourced primarily by collisionless stellar or compact-object components.
  \item Gas thermalisation suppresses coherent scale-access momentum.
  \item Projection residuals couple to identity persistence rather than raw baryonic mass.
  \item The IKK residual explains galaxy rotation curves but not cluster lensing,
        requiring ordinary dark matter or another sector at cluster scales.
\end{itemize}
The fourth option must remain allowed. Better an honest hybrid model than a grand
unified failure balloon.

%% ================================================================
\section{Failure Conditions}
%% ================================================================

EVM-V must define failure conditions before fitting the data. Otherwise the analysis
becomes astrology with better typography.

\begin{definition}[Weakened]
The model is \emph{weakened} if $\LamPi$, $\pzero$, $\rc$ vary by more than an order
of magnitude across galaxies without correlation to physical observables.
\end{definition}

\begin{definition}[Strongly weakened]
The model is \emph{strongly weakened} if $\mathrm{BIC}_{\mathrm{IKK}} >
\mathrm{BIC}_{\mathrm{MOND}}$ and $\mathrm{BIC}_{\mathrm{IKK}} >
\mathrm{BIC}_{\mathrm{NFW}}$ for most galaxies despite having comparable or greater
parameter freedom.
\end{definition}

\begin{definition}[Rotation-model falsification]
The model is \emph{falsified as a useful first-order rotation model} if $\vmod(r)$
cannot match SPARC rotation curves at the level achieved by standard halo models even
with more parameters.
\end{definition}

\begin{definition}[Descriptive curve fitting]
The model is \emph{reclassified as descriptive curve fitting} if $\pu(r)$ must be
independently hand-shaped for every galaxy with no population-level structure.
\end{definition}

\begin{definition}[Theoretically weakened]
The model is \emph{theoretically weakened} if $\GPi(r)$ must be manually adjusted to
avoid solar-system and compact-system constraints rather than emerging from the scale
interval structure of IKK IV.
\end{definition}

\begin{definition}[Cluster-incomplete]
The model is \emph{cluster-incomplete} if it cannot explain lensing offsets in merging
clusters without invoking a completely separate dark matter sector with no principled
boundary from the galactic sector.
\end{definition}

These are not admissions of defeat. They are guardrails. Without them, the theory
becomes immortal in the worst possible way: impossible to kill because it never actually
stands anywhere.

%% ================================================================
\section{Proposed Figures and Tables}
%% ================================================================

\paragraph{Figure~1 --- EVM-V inference chain.}
\[
    \vobs(r) \;\to\; \aobs(r) \;\to\; \abar(r) \;\to\; \aPi(r)
    \;\to\; \text{IKK\,/\,MOND\,/\,halo comparison.}
\]

\paragraph{Figure~2 --- Projection-residual interpretation.}
Show $\Itrue \to \Pab \to \Osp^{(a)}$ with a residual branch
$\Rresid = \Itrue - \Ihat$.

\paragraph{Figure~3 --- Rotation curve decomposition.}
Plot $\vobs$, $\vbar$, $\vPi$, $\vmod$ for representative SPARC galaxies at a range
of stellar masses and surface brightnesses.

\paragraph{Figure~4 --- RAR comparison.}
Plot $\aobs$ vs.\ $\abar$ with curves for the empirical RAR, MOND, and the IKK
residual model overlaid.

\paragraph{Figure~5 --- Parameter stability.}
Plot fitted $\LamPi$, $\pzero$, and $\rc$ across all galaxies with bootstrap confidence
intervals, ordered by stellar mass.

\bigskip
\begin{longtable}{@{}p{2.8cm}p{2.2cm}p{2.5cm}p{3.0cm}p{3.0cm}@{}}
\caption{Model comparison summary\label{tab:models}}\\
\toprule
\textbf{Model} & \textbf{Free params} & \textbf{Data inputs}
  & \textbf{Main strength} & \textbf{Main failure mode}\\
\midrule
\endfirsthead
\toprule
\textbf{Model} & \textbf{Free params} & \textbf{Data inputs}
  & \textbf{Main strength} & \textbf{Main failure mode}\\
\midrule
\endhead
\midrule\multicolumn{5}{r}{\small(continued\ldots)}\\
\endfoot
\bottomrule
\endlastfoot
Baryons only  & 0--2  & SPARC mass model     & Minimal                    & Fails outer curves \\[4pt]
MOND/RAR      & 1--2  & $\abar$              & Tight empirical law        & Cluster/cosmology stress \\[4pt]
NFW halo      & 2+    & Halo profile         & Cosmology-compatible       & Halo--baryon tuning \\[4pt]
IKK residual  & 4--6  & Baryons + $\pu$ proxy & Ontological residual map  & Proxy degeneracy \\
\end{longtable}

\bigskip
\begin{longtable}{@{}p{1.2cm}p{3.5cm}p{3.5cm}p{4.0cm}@{}}
\caption{Hidden-momentum proxy hierarchy\label{tab:proxies}}\\
\toprule
\textbf{Tier} & \textbf{Proxy} & \textbf{Risk} & \textbf{Recommended use}\\
\midrule
\endfirsthead
\toprule
\textbf{Tier} & \textbf{Proxy} & \textbf{Risk} & \textbf{Recommended use}\\
\midrule
\endhead
\bottomrule
\endlastfoot
I   & $\Mstar\vcirc$               & Circularity                      & Baseline only \\[4pt]
II  & $\Sbar\vbar$                  & Surface-density assumptions       & First serious model \\[4pt]
III & $\rhobar\vbar$                & Vertical structure assumptions    & Advanced \\[4pt]
IV  & Morphology-conditioned        & Overfitting                       & Only after Tier II/III \\
\end{longtable}

%% ================================================================
\section{Minimal Reproducible Computational Workflow}
%% ================================================================

A first implementation should produce the following files per galaxy:
\begin{verbatim}
galaxy_id/
  input/
    sparc_rotation_curve.csv
    sparc_mass_model.csv
  derived/
    acceleration_table.csv
    residual_table.csv
  fits/
    baryon_only_fit.json
    mond_fit.json
    nfw_fit.json
    ikk_projection_fit.json
  figures/
    rotation_curve_decomposition.png
    acceleration_relation.png
    residual_profile.png
  report/
    galaxy_summary.md
\end{verbatim}

Minimum table schema (one row per resolved radius):
\begin{verbatim}
r_kpc, v_obs_kms, sigma_v_kms,
v_gas_kms, v_disk_kms, v_bulge_kms, v_bar_kms,
a_obs, a_bar, a_residual,
p_u_proxy, v_model_ikk, v_model_mond, v_model_nfw
\end{verbatim}

Population-level outputs:
\begin{verbatim}
population_parameter_stability.csv
model_comparison_AIC_BIC.csv
rar_residuals.csv
morphology_correlation_summary.csv
failure_flags.csv
\end{verbatim}

The most important file is \texttt{failure\_flags.csv}, because apparently someone has
to protect the theory from its own optimism.

%% ================================================================
\section{Discussion}
%% ================================================================

EVM-V does not prove the IKK interpretation. Its contribution is methodological. It
converts the \emph{Looking Up} framework from a theoretical proposal into a structured
observational programme.

\subsection{Outcome classification}

\paragraph{Strongest outcome.}
The model fits SPARC rotation curves competitively, $\LamPi$, $\pzero$, and $\rc$
remain approximately stable (variation within a factor of 3 across the sample), the
RAR scatter is preserved or reduced relative to the MOND baseline, and the cluster
transition can be located at a principled scale boundary. In this outcome, the
projection-residual interpretation earns the right to be taken seriously as a candidate
framework alongside MOND and $\Lambda$CDM. It does not win by default, but it passes
the minimum bar for scientific credibility.

\paragraph{Intermediate outcome.}
The model fits most SPARC galaxies acceptably, but $\LamPi$, $\pzero$, $\rc$ show
structured variation correlated with galaxy mass, gas fraction, or surface brightness.
In this case the model is a useful phenomenological classifier: it maps different galaxy
populations onto different projection regimes and predicts where the residual is large,
small, or anomalous. This outcome is scientifically interesting but requires an
explanation of why projection efficiency varies with observable morphological properties.
It establishes population-level phenomenology that may or may not reflect a deeper
mechanism.

\paragraph{Weak outcome.}
Parameters vary by more than an order of magnitude without physical correlation.
Fitting is acceptable for individual galaxies but requires galaxy-by-galaxy $\pu$
tuning equivalent to choosing a halo profile. In this case the IKK residual model has
not added anything over standard halo fitting. It has only added vocabulary.

\paragraph{Falsification outcome.}
The model cannot reproduce SPARC rotation curves at the level of NFW halo fits even
with more parameters, or BIC is substantially worse despite equal parameter freedom.
This outcome would demote EVM-V from a projection-access programme to a falsified
application. It would not falsify the abstract IKK projection framework --- which is
structural, not tied to any one dynamical application --- but it would sever the
\emph{Looking Up} bridge from the formal theory to astrophysical data.

\subsection{Relationship to MOND}

The projection-residual framework makes a different ontological claim than MOND, but
it must reproduce MOND-like phenomenology in the regimes where MOND succeeds. The
natural consistency requirement is that the outer-disk flat-rotation limit of the IKK
model recovers the deep-MOND scaling $\aMOND \approx \sqrt{a_0\abar}$. Whether this
requires fine-tuning of $\LamPi$ and $\pzero$ to match $a_0$, or whether it emerges
naturally from the projection-access structure, is a key diagnostic. If the former, the
model is a reparametrisation of MOND rather than an explanation of it.

\subsection{Relationship to dark matter}

Halo fits introduce spatially distributed mass profiles. Projection residuals introduce
a scale-access coupling that does not require extra mass. The two frameworks are
phenomenologically distinguishable only where they predict different radial acceleration
profiles, RAR scatter patterns, or cluster-scale behaviour. EVM-V is specifically
designed to identify those distinguishing observables. If no such observable can be
found at the SPARC level, the two frameworks are empirically degenerate at galactic
rotation-curve scales, and the competition must move to cluster lensing, galaxy
formation statistics, or direct detection programmes.

\subsection{Fifth-force interpretation and theoretical coherence}

The term $\aPi(r,\usc)$ is described throughout as a ``fifth-force-like'' contribution.
This language is deliberately cautious. The IKK framework does not introduce a new
force in the gauge-theory sense. It introduces a projection-access coupling that
produces an effective force at scales where the coupling is non-negligible. The
distinction matters: a gauge force requires a mediator particle and a conserved charge
detectable in particle physics. A projection coupling requires neither, but it also
provides fewer experimental handles.

The theoretical coherence requirement is that $\aPi(r,\usc)$ must connect smoothly to
the formal projection operator machinery of IKK I--IV. Specifically, $\pu(r)$ must be
expressible as a functional of the projection access efficiency $\etaab$, and the
$\GPi(r)$ gate must be derivable from the scale interval structure of IKK IV rather
than chosen purely for curve-fitting convenience. Achieving that connection is a
later-stage theoretical task; EVM-V opens the empirical space that makes such a
connection worth pursuing.

\subsection{The proxy problem is the central risk}

The honest assessment is that $\pu(r)$ is the weakest element of the current
formulation. It is a theoretical quantity with no direct observable. The four proxy
tiers defined in Section~8 are operational stand-ins with varying degrees of physical
motivation. The Tier~I proxy is almost certainly too close to the observable to be
informative. The Tier~II proxy is the most defensible first candidate. The real test of
the framework is whether Tier~II or~III proxies, applied uniformly, produce stable
parameters across the SPARC sample. If they do not, the model has a proxy problem it
cannot escape without additional theoretical input specifying what physical process
sources $\pu$.

%% ================================================================
\section{Conclusion}
%% ================================================================

EVM-V defines a reproducible observational programme for testing whether galactic
rotation-curve residuals can be mapped onto a projection-access model derived from the
IKK framework. Its primary deliverables are:

\begin{enumerate}
  \item A precise definition of the empirical residual
        $\aPi(r) = \aobs(r) - \abar(r)$ and a reproducible pipeline for computing
        it from SPARC data.
  \item A minimal four-to-six-parameter velocity model with an explicit
        hidden-momentum proxy hierarchy and a regime gate, descended from the
        \emph{Looking Up} construction.
  \item A four-model comparison framework (baryons only, MOND/RAR, NFW halo, IKK
        residual) with AIC/BIC model-selection metrics.
  \item Population-level stability tests for $\LamPi$, $\pzero$, $\rc$ with explicit
        stability thresholds and correlation tests against galaxy observables.
  \item Six named failure conditions that must be evaluated before any positive claim
        is made.
  \item A minimal reproducible workflow schema specifying per-galaxy and
        population-level output files.
\end{enumerate}

The failure conditions are not secondary. They are the primary scientific contribution.
A theory that specifies how it can be shown wrong is a scientific theory. A theory that
does not is a narrative. EVM-V insists on the former.

The expected next step is EVM-VI, which would apply this pipeline to the full SPARC
sample and report results. Whether that paper finds stability, structured variation, or
falsification is genuinely unknown. That is the point.

%% ================================================================
%%  APPENDIX: SYMBOL TABLE
%% ================================================================

\appendix
\section{Symbol Table}

\begin{longtable}{@{}p{4.0cm}p{6.5cm}p{2.0cm}@{}}
\caption{Principal symbols and their definitions\label{tab:symbols}}\\
\toprule
\textbf{Symbol} & \textbf{Meaning} & \textbf{Section}\\
\midrule
\endfirsthead
\toprule
\textbf{Symbol} & \textbf{Meaning} & \textbf{Section}\\
\midrule
\endhead
\midrule\multicolumn{3}{r}{\small(continued\ldots)}\\
\endfoot
\bottomrule
\endlastfoot
%% IKK framework
$\Ia$           & Source identity at scale $S_a$                     & \S3.1 \\[2pt]
$\Itrue$        & Primitive identity (fraktur)                        & \S3.1 \\[2pt]
$\Pab$          & Projection operator $S_a \to S_b$                   & \S3.1 \\[2pt]
$\Osp^{(a)}$    & Observable at scale $S_b$ of source $a$             & \S3.1 \\[2pt]
$\Rresid$       & Projection residual $\Itrue - \Ihat$                & \S3.3 \\[6pt]
%% Kinematics
$r$             & Galactic radius                                     & \S5.1 \\[2pt]
$\vobs(r)$      & Observed circular velocity                          & \S5.1 \\[2pt]
$\sigma_v(r)$   & Velocity measurement uncertainty                    & \S5.1 \\[2pt]
$\vbar(r)$      & Baryonic model velocity                             & \S5.2 \\[2pt]
$\vgas,\vdisk,\vbulge$ & Gas, disk, bulge velocity components          & \S5.2 \\[2pt]
$\Upsdisk,\Upsbulge$   & Stellar mass-to-light ratios                  & \S5.2 \\[6pt]
%% Accelerations
$\aobs(r)$      & Observed centripetal acceleration                   & \S5.1 \\[2pt]
$\abar(r)$      & Baryonic Newtonian acceleration                     & \S5.2 \\[2pt]
$\aPi(r,\usc)$  & Projection-residual acceleration                    & \S3.3 \\[2pt]
$a_0$           & MOND characteristic acceleration scale              & \S6   \\[2pt]
$\aMOND(r)$     & MOND acceleration                                   & \S6   \\[6pt]
%% Scale coordinate
$\usc = \ln(\ell/\ellz)$ & Logarithmic scale coordinate              & \S3.2 \\[2pt]
$\ellz$         & Reference scale                                     & \S3.2 \\[2pt]
$\pw$           & Legacy hidden-axis momentum (deprecated)            & \S3.2 \\[2pt]
$\pu(r)$        & Modernised scale-access momentum proxy              & \S3.2 \\[2pt]
$\etaPi(r,\usc)$ & Projection-access efficiency                       & \S7.3 \\[2pt]
$\ascale(r)$    & Scale-residual acceleration kernel                  & \S7.3 \\[6pt]
%% Model parameters
$\LamPi$        & Projection-residual amplitude                       & \S7.1 \\[2pt]
$\pzero$        & Hidden-momentum normalisation scale                 & \S7.1 \\[2pt]
$\rc$           & Core suppression radius                             & \S7.1 \\[2pt]
$\rzero$        & Gate activation radius                              & \S7.1 \\[2pt]
$k$             & Gate steepness                                      & \S7.1 \\[2pt]
$\alpha_u$      & Proxy calibration coefficient                       & \S8   \\[2pt]
$\GPi(r)$       & Sigmoid activation gate                             & \S7.1 \\[6pt]
%% Galaxy observables
$\Sbar(r)$      & Baryonic surface density                            & \S8   \\[2pt]
$\rhobar(r)$    & Baryonic volume density                             & \S8   \\[2pt]
$\Mstar$        & Stellar mass                                        & \S8   \\[2pt]
$\fgas$         & Gas fraction                                        & \S8   \\[2pt]
$\Rd$           & Disk scale length                                   & \S8   \\[6pt]
%% Statistics
$\chi^2,\chi^2_\nu$ & Chi-squared, reduced chi-squared               & \S10.3 \\[2pt]
AIC, BIC        & Akaike, Bayesian information criteria               & \S10.5 \\[2pt]
$S_\theta$      & Population stability statistic                      & \S11.1 \\[2pt]
$\sigma_{\mathrm{RAR}}$ & RAR scatter                                & \S11.3 \\
\end{longtable}

\end{document}
